\subsection{Modules and Submodules}

Throughout this subsection, $R$ is a nonzero commutative ring with identity $1_R$.\\
Modules replace field scalars by ring scalars.\\
Nonzero scalars may therefore be noninvertible or annihilate nonzero vectors.

\begin{definition} \label{def:module-over-ring} \hlt{$R$-Module}\\
An \hlt{$R$-module} is an abelian group $(M,+)$ equipped with an action $R\times M\to M$, $(r,u)\mapsto ru$.\\
For $(r,s,u,v)\in R^2\times M^2$, the following identities hold:
\begin{enumerate}[label=(\roman*)]
\setlength{\itemsep}{0pt}
\item $r(u+v)=ru+rv$.
\item $(r+s)u=ru+su$.
\item $(rs)u=r(su)$.
\item $1_Ru=u$.
\end{enumerate}
The ring $R$ is the \hlt{base ring} of $M$; denote the additive identity of $M$ by $0_M$.
\end{definition}

\begin{example} \label{ex:basic-module-models} \hlt{Basic Module Models}\\
Fix $m,n\in\N$.
\begin{enumerate}[label=(\roman*)]
\setlength{\itemsep}{0pt}
\item $R^n$ is an $R$-module under componentwise addition and scalar multiplication.
\item $\operatorname{Mat}(m,n,R)$ is an $R$-module under entrywise operations.
\item $R$ is an $R$-module over itself by ring multiplication.
\item Every vector space over a field $\mathbb{F}$ is an $\mathbb{F}$-module.
\end{enumerate}
The module axioms hold componentwise in (i) and entrywise in (ii). They are ring laws in (iii) and vector-space axioms in (iv).
\end{example}

\begin{perspective} \label{per:base-ring-controls-modules} \hlt{The Base Ring Controls the Geometry}\\
The module axioms separate the additive group from the scalar ring.\\
Over a field, scalar division supports the usual theory of bases, dimension, and complements.\\
Over a general ring, nonunits and zero divisors can obstruct division and cancellation.\\
Consequently, bases, dimension, and complements need not retain their vector-space behavior.
\end{perspective}

\begin{definition} \label{def:submodule} \hlt{Submodule}\\
Let $M$ be an $R$-module. A \hlt{submodule} $N\leq M$ is a subset $N\subseteq M$ that is an $R$-module under the restricted operations.
\end{definition}

\begin{theorem} \label{thm:submodule-test} \hlt{Submodule Test}\\
Let $M$ be an $R$-module and $\emptyset\neq N\subseteq M$. Then $N\leq M\Longleftrightarrow(\forall\,(r,s,u,v)\in R^2\times N^2)\,ru+sv\in N$.
\begin{proof}
\begin{enumerate}[label=\roman*.]
\setlength{\itemsep}{0pt}
\item[$(\Rightarrow)$] Assume $N\leq M$. Fix $(r,s,u,v)\in R^2\times N^2$. Then $ru,sv,ru+sv\in N$ by scalar and additive closure.
\item[$(\Leftarrow)$] Assume $(\forall\,(r,s,u,v)\in R^2\times N^2)\,ru+sv\in N$. $(\forall\,x\in M)\,0_Rx=0_M$ by $0_Rx=(0_R+0_R)x=0_Rx+0_Rx$, cancellation in $(M,+)$.\\
$(\forall\,x\in M)\,(-1_R)x=-x$ by $(-1_R)x+x=0_Rx=0_M$. Fix $n\in N$. Then $0_M=0_Rn+0_Rn\in N$ by the assumed closure.\\
$(\forall\,(x,y)\in N^2)\,x+y=1_Rx+1_Ry\in N$ by the assumed closure. $(\forall\,x\in N)\,-x=(-1_R)x+0_Rx\in N$ by the assumed closure.\\
$(\forall\,(r,x)\in R\times N)\,rx=rx+0_Rx\in N$ by the assumed closure. $N\leq(M,+)$ by the preceding zero, addition, and inverse relations.\\
The set $N$ is closed under $R$-scalar multiplication by the final relation.\\
$N\leq M$ by Definition \ref{def:submodule}, since the remaining module axioms are inherited from $M$.
\end{enumerate}
\hlt{Structural Commentary}: The reverse direction extracts every submodule operation from two-term closure.\\
\hlt{Key Insight}: Two-term $R$-linear combinations encode all submodule closure laws.
\end{proof}
\end{theorem}

\begin{proposition} \label{prop:submodule-intersection-sum} \hlt{Intersections and Sums of Submodules}\\
Let $M$ be an $R$-module and $A,B\leq M$. Define $\mathcal{S}_{A,B}\defeq\{N\leq M:A\cup B\subseteq N\}$.
\begin{enumerate}[label=(\roman*)]
\setlength{\itemsep}{0pt}
\item $A\cap B\leq M$.
\item $A+B\defeq\{a+b:(a,b)\in A\times B\}\leq M$.
\item $A+B=\bigcap\limits_{N\in\mathcal{S}_{A,B}}N$.
\end{enumerate}
\begin{proof}
\begin{enumerate}[label=(\roman*)]
\setlength{\itemsep}{0pt}
\item $0_M\in A\cap B$ by $A,B\leq M$. For $(r,s,x,y)\in R^2\times(A\cap B)^2$, $rx+sy\in A\cap B$ by closure of $A$, $B$; $A\cap B\leq M$ by Theorem \ref{thm:submodule-test}.
\item $0_M=0_M+0_M\in A+B$. Fix $x=a+b$ and $y=a'+b'$ with $(a,a',b,b')\in A^2\times B^2$.\\
For $(r,s)\in R^2$, $rx+sy=(ra+sa')+(rb+sb')\in A+B$ by closure of $A$ and $B$. Hence $A+B\leq M$ by Theorem \ref{thm:submodule-test}.
\item $(\forall\,a\in A)\,a=a+0_M\in A+B$ and $(\forall\,b\in B)\,b=0_M+b\in A+B$. $A\cup B\subseteq A+B$ by these inclusions, so $A+B\in\mathcal{S}_{A,B}$ by (ii).\\
$(\forall\,N\in\mathcal{S}_{A,B})\,A+B\subseteq N$ by additive closure in $N$. Hence $A+B\subseteq\bigcap\limits_{N\in\mathcal{S}_{A,B}}N$ by the definition of intersection.\\
Conversely, $\bigcap\limits_{N\in\mathcal{S}_{A,B}}N\subseteq A+B$ by $A+B\in\mathcal{S}_{A,B}$.
\end{enumerate}
\hlt{Structural Commentary}: The closure test constructs both submodules; universal containment identifies their join.\\
\hlt{Key Insight}: $A+B$ contains exactly the additions forced by $A$ and $B$.
\end{proof}
\end{proposition}

\begin{proposition} \label{prop:submodules-of-regular-module-are-ideals} \hlt{Ideals Are Submodules of the Regular Module}\\
Equip $(R,+)$ with the action $R\times R\to R$, $(r,x)\mapsto rx$. Write ${}_R R$ for this regular module. Then $(\forall\,N\subseteq R)\,[N\leq{}_R R\Longleftrightarrow N\trianglelefteq R]$.
\begin{proof}
The group $(R,+)$ is abelian by Definition \ref{def:ring-commutative-identity-field}. The ring laws make ${}_R R$ an $R$-module by Definition \ref{def:module-over-ring}.
\begin{enumerate}[label=\roman*.]
\setlength{\itemsep}{0pt}
\item[$(\Rightarrow)$] Assume $N\leq{}_R R$. Then $N\neq\emptyset$ and $(\forall\,a,b\in N)\,a-b\in N$ by Definition \ref{def:submodule}.\\
$(\forall\,(r,n)\in R\times N)\,rn,nr\in N$ by scalar closure and commutativity of $R$. Hence $N\trianglelefteq R$ by Definition \ref{def:ideal-and-generated-ideal}.
\item[$(\Leftarrow)$] Assume $N\trianglelefteq R$. Then $N\leq(R,+)$ by Proposition \ref{prop:subgroup-test} and Definition \ref{def:ideal-and-generated-ideal}.\\
For $(r,s,u,v)\in R^2\times N^2$, $ru,sv\in N$ by ideal absorption and $ru+sv\in N$ by additive closure. Hence $N\leq{}_R R$ by Theorem \ref{thm:submodule-test}.
\end{enumerate}
\hlt{Structural Commentary}: The regular action identifies scalar closure with ideal absorption; commutativity adds right absorption.\\
\hlt{Key Insight}: Submodules of ${}_R R$ and ideals share the same closure laws.
\end{proof}
\end{proposition}
